\chapter{Étape 1 contre Étape 2 : discussion critique}
\label{ch-disc}

\section{Ce que la simulation confirme et ce qu'elle corrige}

\begin{table}[H]
\centering
\caption{Étape 1 et Étape 2 comparées.}
\label{tab-compare}
\begin{tabularx}{\textwidth}{@{}lLLl@{}}
\toprule
\thead{Grandeur} & \thead{Étape 1 (analytique)} & \thead{Étape 2 (Atoll)} & \thead{Verdict} \\
\midrule
Nombre de sites & 21 & 21, positions différentes ($-$R02, +R11, +R12) & confirmé \\
Couverture urbaine indoor & objectif \SI{95}{\percent} & \SI{99,7}{\percent} & atteint \\
Couverture rurale outdoor & objectif \SI{95}{\percent} & \SI{94,3}{\percent} de la surface, \SI{97,9}{\percent} de la population & $-0{,}7$ pt \\
Liaison limitante & montante, \SI{2}{\deci\bel} sous la descendante & DCFI UL insensible à la charge : limite de puissance du terminal & confirmé \\
Capacité & 5 sites suffiraient dans le centre & aucune saturation, même à \num{1000} utilisateurs & confirmé \\
Interférence ($K = 1$) & \SI{8,4}{\percent} de la cellule idéale sous \SI{-1}{\deci\bel} & \SI{21}{\percent} de la ville sous \SI{-1}{\deci\bel} avec la charge simulée & optimiste \\
Forme des cellules & cercles réguliers de \SI{0,64}{\kilo\metre} et \SI{2,06}{\kilo\metre} & ombres du relief au sud-ouest et au pied est & relief ignoré \\
\bottomrule
\end{tabularx}
\end{table}

La simulation confirme le nombre de sites, la liaison limitante et la conclusion sur la capacité : le calcul analytique a donné le bon ordre de grandeur. Quatre points changent.

\begin{itemize}
  \item Les positions. La formule de Hata ne voit pas le relief ; Atoll, avec la diffraction, calcule l'obstruction de la crête. Deux sites ont été ajoutés à l'intérieur des ombres du massif, et un site de la grille régulière s'est révélé inutile.
  \item L'interférence. L'hexagone idéal suppose des voisins réguliers sur un terrain plat. Le réseau réel a des sites perchés qui débordent loin dans la plaine, et Atoll applique la marge d'effet de masque aussi au C/(I+N). Le DL en ville reste ainsi à \SI{79}{\percent} de la surface au lieu de \SI{95}{\percent} ; sur ce point, notre calcul de l'Étape 1 était optimiste.
  \item La charge. Les prédictions par défaut supposent toutes les cellules chargées à \SI{100}{\percent}, alors que la charge simulée est de \SIrange{4}{31}{\percent}. Une même carte peut paraître très mauvaise ou acceptable selon cette hypothèse, qu'il faut donc toujours préciser.
  \item Le bord de cellule en UL. Dans l'Étape 1, le seuil de \SI{-120}{\dBm} et le débit UL de \SI{128}{\kbps} décrivent le même point. Dans Atoll, la surface couverte en outdoor (\SI{73,8}{\kilo\metre\squared}) dépasse la surface où l'UL atteint \SI{128}{\kbps} (\SI{48,6}{\kilo\metre\squared}). Les surcoûts et le taux d'erreur appliqués par Atoll pourraient exiger un troisième RB en bord de cellule, soit $10\log_{10}(3/2) \approx \SI{1,8}{\deci\bel}$ de moins par RB ; nous ne l'avons pas vérifié.
\end{itemize}

\section{Les notions du cours dans le projet}
\label{sec-cours}

Le \cref{tab-cours} reprend les notions du cours~\cite{afif2026cours} utilisées dans ce rapport, la façon dont nous les avons appliquées au LTE et le résultat obtenu.

\begin{table}[H]
\centering
\caption{Notions du cours et leur application.}
\label{tab-cours}
\small
\begin{tabularx}{\textwidth}{@{}>{\raggedright\arraybackslash}p{0.23\textwidth}>{\raggedright\arraybackslash}p{0.26\textwidth}L@{}}
\toprule
\thead{Notion du cours} & \thead{Formule ou règle} & \thead{Application et résultat} \\
\midrule
Objectif selon le contexte & urbain : capacité ; rural : couverture & trois zones ; la ville reste limitée par la couverture aujourd'hui (11 sites contre 5) \\
Bilan de liaison équilibré & bord de cellule atteint en UL et en DL & deux bilans ; UL limitante de \SI{2}{\deci\bel}, $R = \SI{0,64}{\kilo\metre}$ et \SI{2,06}{\kilo\metre} \\
Marge de fading & 10 à \SI{15}{\deci\bel} (FDMA) & \SI{0}{\deci\bel} en LTE, déjà dans les seuils des bearers ; sinon 3,7 fois plus de sites \\
Site tri-sectoriel & un site couvre trois cellules & 21 sites, 63 cellules ; $A_{\text{site}} = 1{,}95\,R^2$ \\
Taille du motif & $K = i^2 + ij + j^2$ & $K = 1$ avec $(i, j) = (1, 0)$ \\
Distance de réutilisation & $D = \sqrt{3K}\,R$, $q = D/R$ & $q = 1{,}73$ pour $K = 1$ \\
C/I au premier anneau & $C/I = (q - 1)^{\alpha} / 6$ & $\alpha = 3{,}52$ (pente de Hata) ; \SI{-12,6}{\deci\bel} au pire cas, \SI{8,4}{\percent} de la cellule sous \SI{-1}{\deci\bel} \\
Seuils de C/I & GSM \SI{9}{\deci\bel} : $K = 7$ ou 9 ; \SI{3}{\deci\bel} : $K = 3$ ou 4 & retrouvés par le calcul ; LTE à \SI{-1}{\deci\bel}, d'où $K = 1$ \\
Rendement d'un canal & $\eta$ = canaux de trafic / canaux totaux & GSM 0,75 ; LTE 0,71 en DL et 0,86 en UL \\
Capacité & $n\,\eta / (K \cdot B \cdot \pi R^2)$ & leviers utilisés : $K = 1$ et taille des cellules ; 5 sites suffisent aujourd'hui, 11 dans cinq ans \\
Boucle de planification & prédictions, ajustement, densification & huit actions mesurées, dont deux sites ajoutés \\
\bottomrule
\end{tabularx}
\end{table}

\section{Limites de l'étude}
\label{sec-limites}

\begin{itemize}
  \item Modèle non calibré. Aucune mesure terrain n'a permis d'ajuster le modèle Cost-Hata. L'incertitude de chaque pixel ($\sigma$ de 6 à \SI{8}{\deci\bel}) est bien plus grande que les 0,7 point qui manquent en campagne.
  \item Extrapolation du modèle. Le modèle de Hata est validé à partir de \SI{1}{\kilo\metre} ; le rayon urbain de \SI{0,64}{\kilo\metre} en est une légère extrapolation.
  \item Critère de qualité sévère. La marge d'effet de masque de \SI{86}{\percent} est appliquée aussi au C/(I+N) ; sans elle, la carte de qualité est nettement meilleure.
  \item Ordonnanceur d'Atoll. Même à 10 utilisateurs, \SI{46}{\percent} des utilisateurs HSI n'obtiennent pas \SI{1}{\Mbps} alors qu'ils sont seuls dans leur cellule avec un bon C/I. Nous n'avons pas trouvé la règle interne qui l'explique.
  \item Trafic modélisé. Un habitant vaut un utilisateur potentiel, et le mélange de services, de mobilités et de terminaux est une hypothèse.
  \item Antenne. Le diagramme d'antenne utilisé est défini à \SI{2100}{\mega\hertz} ; la bibliothèque d'Atoll contient un modèle à \SI{1800}{\mega\hertz} (\ang{65}, \SI{17}{\dBi}, tilt électrique de \ang{6}) qui serait plus juste.
  \item Objectifs non tous atteints. La campagne reste à \SI{94,3}{\percent} de la surface (\SI{97,9}{\percent} de la population) et le DL en ville à \SI{79}{\percent} de la surface.
\end{itemize}

\section{Perspectives}

Les réglages d'antenne ont atteint leur limite (action 7 rejetée). Les pistes suivantes portent sur l'architecture du réseau :
\begin{itemize}
  \item la coordination d'interférence entre cellules (ICIC) pour les bords de cellule urbains ;
  \item une couche à \SI{800}{\mega\hertz} sur les sites existants, qui apporterait de la portée en campagne et à l'intérieur des bâtiments, en complément du \SI{1800}{\mega\hertz} qui porte la capacité (\cref{sec-800-1800}) ;
  \item des petites cellules dans le centre quand le trafic aura doublé, au rythme actuel de +\SI{40}{\percent} tous les deux ans ;
  \item une campagne de mesures pour calibrer le modèle de propagation avant tout déploiement réel.
\end{itemize}
